Several limitations remain. First, the dataset is a curated enzyme optimal-pH resource whose measurements and source annotations were inherited from the EpHod release. Zhou and Fang identified disputed pH values and protein-identity assignments in this resource \cite{zhou2026caution}; our audit confirmed that B5LXD7 is present and quantified its influence, but did not independently recheck every flagged record against every original experiment or alter labels. Consequently, extreme-pH errors should not be interpreted as model failures without source-level verification. Second, taxonomy metadata beyond organism labels were not available in the processed file, so taxonomy-held-out evaluation was limited to organism-held-out splitting. Third, the 50\%, 70\%, and 90\% MMseqs2 cluster thresholds were evaluated, but cluster composition and target-distribution differences still complicate direct comparisons across thresholds.

Fourth, the organism-held-out and EC-held-out splits are group-disjoint but not necessarily distribution-matched. Their test sets may differ in target-pH variance and group-size composition from the generated random test set, which complicates direct MAE comparisons. Fifth, subgroup analyses are limited by sample size; results for rare organisms or EC numbers should not be overinterpreted.

Finally, the ESM-2 embeddings were used as frozen representations with classical regressors. The study does not fine-tune protein-language models, does not reproduce any external model's complete training protocol, and does not claim biological mechanisms from embedding-space or error-analysis patterns.
